\subsection{幺半群作用于集合}

定义1. 设M为一个幺半群，其运算用乘法表示，单位元记为e，E为一个集合。一个作用 \(\alpha \mapsto f_{\alpha}\) M在E上的作用称为M在E上的左（相应地，右）作用，如果 \(f_{e} = \mathrm{Id}_{\mathrm{E}}\) 和 \(f_{\alpha \beta} = f_{\alpha} \circ f_{\beta}\) （相应地 \(f_{\alpha \beta} = f_{\beta} \circ f_{\alpha}\) ）对所有 \(\alpha, \beta \in \mathbf{M}\) .

换言之，幺半群M在集合E上的左（相应地，右）作用是从M到幺半群 \(E^{E}\) （分别为 \(E^{E}\) 的相反幺半群）与映射的复合。若与M的作用相对应的作用律记为左（分别为右）乘法，则该作用为左（分别为右）运算这一事实可用公式表示为

\[
e. x = x; \alpha . (\beta . x) = (\alpha \beta). x \quad \text { for } \alpha , \beta \in \mathrm{M} \text { and } x \in \mathrm{E}.\tag{1}
\]

\[
(\text { resp. } x. e = x; (x. \alpha). \beta = x. (\alpha \beta) \quad \text { for } \alpha , \beta \in M \text { and } x \in E).
\]

在这些条件下，也称M在E上左（相应地，右）作用，且相应的作用法则为幺半群M在E上的左（相应地，右）运算律。

设 M 是一个幺半群；带有 M 在 E 上的左（相应地，右）作用的集合 E 称为左（相应地，右）M-集。若映射 \(\alpha \mapsto f_{\alpha}\) 的M到 \(E^{E}\) 是单射。

设 E 为一个集合；其典范作用为 \(E^{E}\) 在E（§3，第1号，例3）上是左作用。

\begin{enumerate}
\def\labelenumi{(\arabic{enumi})}
\setcounter{enumi}{1}
\tightlist
\item
  设 M 是一个幺半群。由 M 上的运算导出的 M 对自身的左（相应地，右）作用（§ 3，第 3 号，例 7）是 M 对自身的一个左（相应地，右）运算。在考虑这一运算时，我们说 M 通过左（相应地，右）平移作用于自身。
\end{enumerate}

设 E 是一个左（相应地，右）M-集合， \(M^{0}\) M 的对偶幺半群。在相同作用下，幺半群 \(M^{0}\) 在 E 上右（相应地，左）作用。所得的 \(M^{0}\) -集合称为 M-集合 E 的对偶集合。关于左 M-集合的定义和结论可平移至右 \(M^{0}\) -集，当过渡到相反结构时。

在本段其余部分中，除非另有说明，我们将仅考虑左M-集，并简称为M-集。其作用法则将用左乘法表示。

设E为一个集合。设G为作用于E上的一个群。对于所有 \(\alpha\) 在 G 中，元素 \(E^{E}\) 上的作用，定义为 \(\alpha\) 是 E 的一个置换（ \(\S\) 2，第 3 号，例 2）。给定 G 在 E 上的一个作用，因此等价于给定一个从 G 到 \(S_{E}\) .

依照第3节第3项，我们作出如下定义：

定义2. 设 M 为一个幺半群，且 E 与 E' 为 M-集合。称映射 f 从 E 到 E' 满足，对所有 \(x \in E\) 以及所有 \(\alpha \in M\) , \(f(\alpha \cdot x) = \alpha \cdot f(x)\) 被称为M-集同态（或M-态射，或与M的运算相容的映射）。

M-集合的恒等映射是M-态射。两个M-态射的复合仍是M-态射。一个M-集合到另一个M-集合的映射成为同构的充要条件是：它是一个双射的M-态射，并且其逆映射也是M-态射。

设 \((\mathbf{E}_i)_{i\in \mathbf{I}}\) 设为一族M-集合，并令 \(\mathbf{E}\) 为 \(\mathbf{E}_i\) 。幺半群 M 作用于 \(\mathbf{E}\) 中条件的扩张 \(\alpha .(x_i)_{i\in \mathbf{I}} = (\alpha .x_i)_{i\in \mathbf{I}}\) 和 \(\mathbf{E}\) ，在此作用下，是一个M-集；设 \(\mathbf{E}'\) 是一个M-集；映射 \(f\) 于 \(\mathbf{E}'\) 到 \(\mathbf{E}\) 是M-态射当且仅当 \(\mathbf{pr}_i\circ f\) 是 \(\mathbf{E}'\) 到 \(\mathbf{E}_i\) 为所有 \(i\in \mathbf{I}\) .

的M-态射。设E是一个M-集，F是E在M作用下稳定的子集；在诱导的运算下，F是一个M-集，且典范嵌入 \(F \rightarrow E\) 是M-态射。设E是一个M-集，R是E上与M作用相容的等价关系；商集E/R在商作用下是一个M-集，且典范映射

是幺半群同态， \(E \rightarrow E/R\) 是M-态射。设E是一个M-集，R是E上与M作用相容的等价关系；商集E/R在商作用下是一个M-集，且典范映射

设 \(\phi: \mathbf{M} \to \mathbf{M}'\) -集和 \(\mathbf{E}\) 一个 \(\mathbf{M}\) -集。映射 \(\mathbf{E}'\) 一个 \(\mathbf{M}'\) 满足对所有 \(f\) 于 \(\mathbf{E}\) 到 \(\mathbf{E}'\) 有 \(x \in \mathbf{E}\) 和 \(\alpha \in \mathbf{M}\) ,

\[
f (\alpha . x) = \phi (\alpha). f (x)
\]

称为 \(\phi\) -从E到E'的态射（参见§3，第1号）。

运算律的推广。给定（例如）三个集合 \(F_{1}, F_{2}, F_{3}\) ，排列 \(f_{1}, f_{2}, f_{3}\) 于 \(F_{1}, F_{2}, F_{3}\) 分别，以及基集上的一个阶梯F \(F_{1}, F_{2}, F_{3}\) （集合论，IV，§ 1，第 1 号），我们可以沿着分层结构 F 的构造逐步定义 F 的一个置换，称为 \(f_{1}, f_{2}, f_{3}\) 至F（集合论，IV，§1，第2号）；我们将其记为 \(\phi_{F}(f_{1}, f_{2}, f_{3})\) .

那么设 G 为一个群，且 \(h_{i}\) G 到对称群的一个同态， \(F_{i}\) (i = 1, 2, 3)，换言之，G 在 \(F_{i}\) 该映射 \(x \mapsto x_{F} = \phi_{F}(h_{1}(x), h_{2}(x), h_{3}(x))\) 是G到 \(S_{F}\) ，换言之，即G在F上的一个作用，称为 \(h_{1}, h_{2}, h_{3}\) 到F。设P是F的一个子集，使得对于所有 \(x \in G\) , \(x_{F}(P) = P\) ；令 \(x_{P}\) 为 \(x_{F}\) 到

P；则映射 \(x \mapsto x_{P}\) 是G在P上的一个运算，也称为 \(h_{1}, h_{2}, h_{3}\) 到P。

例如，设 \(\mathbf{K}\) 和 \(\mathbf{L}\) 是上的两个层；取 \(\mathbf{F}_1, \mathbf{F}_2, \mathbf{F}_3\) 为 \(\mathbf{F}\) 的子集构成的集合，取 \(\mathbf{K} \times \mathbf{L}\) 和 \(\mathbf{P}\) 为 \(\mathbf{K}\) 到 \(\mathbf{L}\) 的映射构成的集合，与它们的图等同。如果 \(w \in \mathbf{P}\) 和 \(x \in \mathbf{G}\) , \(x_{\mathbf{P}}(w)\) 是映射 \(k \mapsto x_{\mathbf{L}}(w(x_{\mathbf{K}}^{-1}(k)))\) 于 \(\mathbf{K}\) 到 \(\mathbf{L}\) .

